% \section{Standard Maps between Cells and Spheres}
% \begin{definition}
% \begin{enumerate}
%     \item The standard $n-$ball
%     \[\B^n = \{x \in \R^n: \norm{x} \leq 1\}.\]
    
%     \item The standard $n-$phere
%     \[\mathbb{S}^n= \partial(\B^{n+1}) = \{x \in \R^{n+1}: \norm{x} =1\}.\]

%     \item The standard $n-$cell is a topological space that is homeomorphic to the open unit ball 
%     \[\Int(\B^n) = \{x \in \R^n: \norm{x} < 1\}.\]
% \end{enumerate}
% \end{definition}


% \begin{proposition}[Stereographic projection]
%     For $n \geq 0$, denote 
%     \[N=(0,\ldots,0,1),~S=(0,\ldots,0,-1)\in \mathbb{S}^n\]
%     called the north pole and the south pole of $\mathbb{S}^n$, respectively. Then $\mathbb{S}^n\setminus\{S\}$ and $\mathbb{S}^n\setminus\{N\}$ are homeomorphic to $\R^n$.
% \end{proposition}

% \begin{proof}
%     Consider the continuous maps
%     \[
%         f:\mathbb{S}^n\setminus\{N\}\to \R^n,~
%         x=(a,t) \in \mathbb{S}^n\setminus\{N\} \cap \R^n \times \R\mapsto \dfrac{1}{1-t}a,
%     \]
%     \[
%         g:\R^n\to \mathbb{S}^n\setminus\{N\}, \quad 
%         y \mapsto (\phi(y)\cdot y,1-\phi(y)),
%     \]
%     where $\phi(y)=\dfrac{2}{1+\|y\|^2}$.

%     For every $x=(a,t)\in \mathbb{S}^n$, we compute
%     \[
%         \phi(f(x))=\phi\left(\dfrac{1}{1-t}a\right)=\dfrac{2}{1+\dfrac{1}{(1-t)^2}\norm{a}^2}
%         =1-t,
%     \]
%     hence
%     \begin{align*}
%         (g\circ f)(x)
%         &=g\left(\dfrac{1}{1-t}a\right)\\
%         &=\left(\phi(f(x))\cdot f(x),1-\phi(f(x))\right)\\
%         &=(a,t)=x.
%     \end{align*}

%     For every $y\in \R^n$, we have
%     \begin{align*}
%         (f\circ g)(y)
%         &=f(\phi(y)\cdot y,1-\phi(y))\\
%         &= \dfrac{1}{1-(1-\phi(y))}\left(\phi(y)\cdot y\right)\\
%         &=y.
%     \end{align*}

%     Therefore $f$ is a continuous bijection with continuous inverse $f^{-1}=g$, i.e.\ $f$ is a homeomorphism between $\mathbb{S}^n\setminus\{N\}$ and $\R^n$.

%     On the other hand, there is also a homeomorphism.
%     \[
%         \mathbb{S}^n\setminus\{N\}\to \mathbb{S}^n\setminus\{S\}, \quad 
%         (a,t)\mapsto (a,-t).
%     \]
% \end{proof}

% \begin{proposition}
%     The open unit ball $\Int(\B^n)$ 
%     is homeomorphic to $\R^n$.
% \end{proposition}

% \begin{proof}
%     We have constructed explicit continuous maps
%     \[
%         \varphi:\Int(\B^n)\to \R^n, \quad x \mapsto \dfrac{x}{1-\norm{x}}, 
%     \]
%     and
%     \[
%         \psi:\R^n\to \Int(\B^n), \quad y \mapsto \dfrac{y}{1+\|y\|}.
%     \]
%     Both $\varphi$ and $\psi$ are continuous and
%     $\psi\circ \varphi =\Id_{\Int(\B^n)},~ \varphi\circ \psi=\Id_{\R^n}$. 
    
%     Hence $\varphi$ is a homeomorphism
%     \[
%         \Int(\B^n) \approx \R^n
%     \]
    
% \end{proof}

% \begin{remark}
%     \begin{enumerate}
%         \item From propositions 1.2 and 1.3, we obtain the homeomorphisms
%         \[\Int(\B^n) \xlongrightarrow[\varphi]{\approx}\R^n \xlongrightarrow[g]{\approx} \mathbb{S}^n\setminus\{N\}.\]
    
%         \item Since 
%         \[\B^n = \Int(\B^n) \sqcup \partial(\B^n)=\Int(\B^n) \sqcup \mathbb{S}^{n-1}\] and 
%         \[\mathbb{S}^n = \mathbb{S}^n\setminus\{N\} \sqcup \{N\}.\]
%         Naturally, we want to define a map
%         \[\pi: \B^n \to \mathbb{S}^n, \]
%         such that 
%         \[\pi(\Int(\B^n)) = \mathbb{S}^{n}\setminus \{N\}= (g \circ \varphi)(\Int(\B^n)),\quad  \pi(\mathbb{S}^{n-1}) = \{N\}.\]
%         With $x \in \Int(\B^n),~y = \varphi(x) = \dfrac{x}{1-\norm{x}}$, so 
%         \begin{align*}
%             \phi(y)&=\dfrac{2}{1+\norm{y}^2}\\
%             &=\dfrac{2}{1+\dfrac{\norm{x}^2}{(1-\norm{x})^2}} \\
%             &= \dfrac{2(1-\norm{x})^2}{\norm{x}^2 + (1-\norm{x})^2}
%         \end{align*}
%         \begin{align*}
%         (g\circ \varphi)(x) &= g\left(y\right)\\
%             &= \left(\phi(y)\cdot y, 1-\phi(y)\right)\\
%             &=\dfrac{1}{\norm{x}^2+(1-\norm{x})^2}\left(2(1-\norm{x})x, 2\norm{x}-1\right) \in \mathbb{S}^n.
%         \end{align*}.
%         \end{enumerate}
%     Moreover, with $z = (a,t) \in \mathbb{S}^n\setminus\{N\} \cap \R^n \times \R$
%     \begin{align*}
%         (g\circ \varphi)^{-1}(z)
%         &=(\psi\circ f)(z)\\
%         &=\psi\left(\dfrac{1}{1-t}a\right)\\
%         &= \dfrac{a}{(1-t)+\norm{a}} \in \Int(\B^n).
%     \end{align*}
% \end{remark}
% Hence, we have a definition, also a proposition following:
% \begin{proposition}
%     The map 
%     \[\pi:(\B^n, \mathbb{S}^{n-1})\to (\mathbb{S}^n,N),~x \mapsto \dfrac{1}{\norm{x}^2+(1-\norm{x})^2}\left(2(1-\norm{x})x, 2\norm{x}-1\right)\]
%     induces a homeomorphism \[\overline{\pi}:\B^n/\mathbb{S}^{n-1} \approx \mathbb{S}^n.\]
    
%     Moreover, $\pi \mid_{\Int(\B^n)}: \Int(\B^n) \approx \mathbb{S}^n\setminus\{N\}$ is a homeomorphism and the inverse of $\pi \mid_{\Int(\B^n)}$ is
%     \[\rho = \psi \circ f: \mathbb{S}^n\setminus\{N\} \to \Int(\B^n),~\]
% \end{proposition}

% \begin{remark}
%     $\B^n/\mathbb{S}^{n-1} = \B^n/\sim$, where $x\sim y$ if $x = y$ or $x,y\in \mathbb{S}^{n-1}$. 
% \end{remark}

% Now, we will show that simplices, cubes and their products are homeomorphic with balls.

% \begin{lemma}[Normalization]
%     Let $B\subset \R^n$ be a closed $n$--ball of radius $r>0$ centered at $c\in\R^n$. 
%     Then there exists an affine homeomorphism 
%     \[
%         T:\R^n \to \R^n
%     \]
%     such that 
%     \[
%         T(B) = \B^n \quad \text{and} \quad T(\partial B) = \partial \B^n.
%     \]
%     Consequently, 
%     \[
%         (B,\partial B)\approx (\B^n,\partial \B^n).
%     \]
% \end{lemma}

% \begin{proof}
%     Define 
%     \[
%         T:\R^n \to \R^n, \qquad T(x) = \dfrac{1}{r}(x-c).
%     \]
%     This is the composition of a translation by $-c$ and a dilation by factor $1/r$, 
%     hence a homeomorphism of $\R^n$ onto itself. 

%     For $x\in B$ we have $\|x-c\|\le r$, so 
%     \[
%         \|T(x)\| = \left\|\dfrac{x-c}{r}\right\|\le 1,
%     \]
%     i.e.\ $T(x)\in \B^n$. Conversely, if $y\in \B^n$, then $\|y\|\le 1$ and 
%     $x = r y + c\in B$, so $T(B)=\B^n$. The same argument shows 
%     $T(\partial B)=\partial \B^n$. 

%     Therefore $T$ restricts to a homeomorphism of pairs 
%     $(B,\partial B)\approx (\B^n,\partial \B^n)$.
% \end{proof}

% \begin{proposition}
%     If $K$ is a compact convex subset of $\R^n$ which contains an $n-$ball $B$ then there is a standard homeomorphism $(B,\partial(B)) \approx (K,\partial(K))$. 
% \end{proposition}
% \begin{proof}
%     Using the lemma 1.7, we only need prove $(\B^n,\partial(\B^n)) \approx (K,\partial(K))$.

%     For each $u\in\mathbb{S}^{n-1}$, put
%     \[
%         \lambda(u) :=\sup\{ t\ge 0 : tu\in K\},
%     \]
%     then 
%     \[1 \leq \lambda(u) < \infty.\] 
%     Indeed, the compactness of $K$ implies there exists  $R>0$ such that $\norm{x}\leq R$ for all $x\in K$, hence $\lambda(u) <\infty$. Moreover, $\lambda(u)\ge 1$, since
%     \[[0,1]\cdot u = \{tu: t\in [0,1]\}\subset \B^n \subset K.\]

%     By the definition, $\lambda(u)u$ is the \textit{furthest point of} $K$ along the ray $\{tu:t\geq 0\}$ and cannot be outside $K$, because all points $tu$ with $t < \lambda(u)$ are in $K$ and $K$ is closed (compact). 
    
%     If $\lambda(u)u \in \Int(K)$, there exists $\varepsilon > 0$ such that $(\lambda(u) + \varepsilon)u \in K$, a contradiction of the definition of $\lambda(u)$. So $\lambda(u)u \in \partial K$ and is the unique intersection of the ray $\{t u:t\ge 0\}$ with $\partial K$.

%     Now, we prove that $\lambda$ is continuous. 
    
%     Let $u_k\to u$ in $\mathbb{S}^{n-1}$ and put $t_k=\lambda(u_k)$.
%     Since $K$ is compact, $(t_k)$ is bounded; pass to a convergent subsequence
%     $t_k\to t$. Because $K$ is closed, $t u=\lim t_k u_k\in K$, hence $t\le \lambda(u)$.
%     For the reverse inequality, fix $\varepsilon>0$. Then
%     $(\lambda(u)-\varepsilon)u \in \Int(K)$; by openness there is a neighborhood
%     $U$ of $u$ with $(\lambda(u)-\varepsilon) v\in K$ for all $v\in U$, hence
%     $\lambda(v)\ge \lambda(u)-\varepsilon$ for $v\in U$. Taking $v=u_k$ close to $u$ gives
%     $\liminf t_k\ge \lambda(u)-\varepsilon$. Since $\varepsilon>0$ is arbitrary,
%     $\liminf t_k\ge \lambda(u)$. Thus $t=\lambda(u)$, so $\lambda$ is continuous.
    
%     Finally, we contruct a homeomorphim between $\B^n$ and $K$. 
    
%     For all $x\in\B^n$, can be written uniquely as $x=ru$ with $u\in\mathbb{S}^{n-1}$ and
%     $r\in[0,1]$. Define
%     \[
%         F:\B^n\longrightarrow K,\qquad x= r u \mapsto r\lambda(u)u.
%     \]
%     This is well-defined and continuous (product of continuous functions on the
%     polar coordinates). It is bijective with inverse
%     \[
%         F^{-1}:K\longrightarrow \B^n,\qquad 
%         F^{-1}(s v)= \dfrac{s}{\lambda(v)}\,v\quad\bigl(v\in\mathbb{S}^{n-1},\ s\in[0,\lambda(v)]\bigr),
%     \]
%     which is also continuous by continuity of $\lambda$. Hence $F$ is a homeomorphism.
    
%     In particular, if $r=1$ then $F(u)=\lambda(u)u\in\partial K$; if $r<1$ then $F(r u)$ lies on the
%     segment between $0$ and $\lambda(u)u$ and hence in $\Int(K)$.
%     Therefore $F(\partial\B^n)=\partial K$ and $F(\Int\,\B^n)=\Int(K)$. Thus $F:(\B^n,\partial\B^n)\approx (K,\partial K)$.
% \end{proof}

% \begin{corollary}
%     The cube \(I^n=[-1,1]^n\subset\R^n\) is homeomorphic to the (closed) unit ball
% \(\B^n=\{x\in\R^n:\|x\|\le 1\}\). Moreover, this homeomorphism sends
% \(\partial I^n\) onto \(\mathbb{S}^{n-1}=\partial(\B^n)\)
% \end{corollary}
% \begin{proof}
%     Since $\norm{x} \leq 1$ implies $|x_j|\le 1$ for each coordinate $x_j$. Thus $I^n$ is a compact convex set of $\R^n$ containing the unit $n$–ball. 
    
%     Proposition 1.8 applies with $K=I^n$ and $B=\B^n$, yielding a standard homeomorphism
%     \[
%         F:(\B^n,\mathbb{S}^{n-1})\xrightarrow{\ \approx\ }(I^n,\partial I^n),~
%         ru \mapsto r\,\lambda(u)\,u 
%         \qquad (u \in \mathbb{S}^{n-1}, r \in [0,1]).
%     \]
% Fix $u\in\mathbb{S}^{n-1}$, for $t\ge 0$,
% \[
% tu\in I^n 
% \iff t u_i \in I,~\forall i
% \iff |t u_i|\le 1,~\forall i
% \iff t\le \dfrac{1}{|u_i|}\ \text{whenever }u_i\neq 0.
% \]
% Hence the set 
% \[\{t\ge 0: tu\in I^n\}=
% \Bigl[\,0,\ \min_{1\le i\le n,\,u_i\neq 0}\dfrac{1}{|u_i|}\,\Bigr],
% \]
% so its supremum is
% \[
% \lambda(u)=\min_{1\le i\le n,\,u_i\neq 0}\dfrac{1}{|u_i|}
% =\dfrac{1}{\max_{1\le i\le n}|u_i|}
% =\dfrac{1}{\|u\|_\infty}.
% \]
% Taking $t=\dfrac{1}{\|u\|_\infty}$ gives 
% $|t u_i|=\dfrac{|u_i|}{\|u\|_\infty}\le 1$ for all $i$, with equality for at least one index attaining the maximum of $|u_i|$. Thus $tu\in\partial I^n$, so the supremum is attained and the intersection with $\partial I^n$ is exactly at $u/\|u\|_\infty$.

% Consequently, its inverse \(F^{-1}:I^n\to\B^n\) is
% \[
% F^{-1}(x)=
% \begin{cases}
% 0, & x=0,\\[4pt]
% \dfrac{\norm{x}_\infty}{\norm{x}}\,x, & x\neq 0,
% \end{cases}
% \]
% which is continuous on \(I^n\) and bijective with continuous inverse.

% Finally, if \(x\in\partial I^n\) then \(\norm{x}_\infty=1\), so
% \(\|F^{-1}(x)\|=\norm{x}_\infty=1\) and hence \(F^{-1}(x)\in\mathbb{S}^{n-1}\).  
% If \(\norm{x}_\infty<1\) (i.e. \(x\in \Int I^n\)), then
% \(\|F^{-1}\|=\norm{x}_\infty<1\), so \(F^{-1}(x)\in \Int\B^n\).
% Therefore \(F^{-1}(\partial I^n)=\mathbb{S}^{n-1}\) and
% \(F^{-1}(\Int I^n)=\Int\B^n\), completing the proof.
% \end{proof}

% \begin{corollary}
%     Let the standard $n$–simplex be
%     \[
%         \Delta_n=\Bigl\{(t_0,\dots,t_n)\in\mathbb{R}^{n+1}_{\ge 0}\ :\ \sum_{i=0}^n t_i=1\Bigr\},
%     \]
%     and let $\partial \Delta_n=\{t\in\Delta_n:\ t_i=0\ \text{for some }i\}$ be its boundary.
    
%     There exists a linear embedding $i:\Delta_n\hookrightarrow\mathbb{R}^n$ such that $K:=\imath (\Delta_n)$ is a compact convex set containing a Euclidean $n$–ball. Consequently, by radial projection there is a \emph{standard} homeomorphism
%     \[
%         \Phi_\Delta:(\Delta_n,\partial \Delta_n)\;\xrightarrow{\ \cong\ }\;(\B^{\,n},\mathbb{S}^{n-1}).
%     \]
% \end{corollary}

% \begin{proof}
%     Define a linear embedding $\imath:\Delta_n\to\mathbb{R}^n$ by sending the vertices as follows:
% for $0\le j\le n-1$ let $v_j=e_j$ (the $j$th standard basis vector of $\mathbb{R}^n$) and let
% \[
%    v_n=-\sum_{j=0}^{n-1} e_j.
% \]
% On barycentric coordinates $t=(t_0,\dots,t_n)\in\Delta_n$ we put
% \[
%    \imath(t)=\sum_{j=0}^n t_j v_j\in\mathbb{R}^n.
% \]
% Then $K:=\imath(\Delta_n)=\operatorname{conv}\{v_0,\dots,v_n\}$ is a compact convex polytope in $\mathbb{R}^n$. We have two claims following
% \begin{itemize}
%     \item $0\in\Int(K)$.
    
% Indeed,
% \[
%    \frac{1}{n+1}\sum_{j=0}^n v_j
%    =\frac{1}{n+1}\Bigl(\sum_{j=0}^{n-1} e_j - \sum_{j=0}^{n-1} e_j\Bigr)=0,
% \]
% so $0$ is a convex combination of the vertices with all coefficients equal to $1/(n+1)>0$; hence $0$ lies in the (relative) interior of $K$.

% \item \textit{$K$ contains a Euclidean $n$–ball around $0$.}

% Since $K$ is compact and $0\notin\partial K$, the continuous function $x\mapsto\|x\|$ attains a positive minimum on the compact set $\partial K$:
% \[
%    r:=\min\{\|x\|:\ x\in\partial K\}>0.
% \]
% Therefore the closed ball $\overline B(0,r)\subset K$.

% By Proposition 1.8 applied to $K$ (with the ball $\overline B(0,r)$ found above), there is a standard homeomorphism
% \[
%    \Psi:(K,\partial K)\xrightarrow{\approx}( \B^{\,n},\mathbb{S}^{n-1})
% \]
% that sends each $y\in\partial K$ to $y/\|y\|$ and extends linearly on the radial segments from $0$.
% \end{itemize}
% Finally, set $\Phi_\Delta:=\Psi\circ \imath:\Delta_n \to \B^{n}$. Since $i$ is a homeomorphism from $\Delta_n$ onto $K$ and $\Psi$ is a homeomorphism from $K$ onto $\B^{\,n}$, $\Phi_\Delta$ is a homeomorphism.

% It remains to check the boundary correspondence. 

% A point $t\in\Delta_n$ lies in $\partial \Delta_n$ iff at least one barycentric coordinate is $0$, say $t_j=0$; equivalently, $i(t)$ lies in the facet $\operatorname{conv}\{v_0,\dots,\widehat{v_j},\dots,v_n\}\subset\partial K$. Thus $i(\partial \Delta_n)=\partial K$, and hence
% \[
%    \Phi_\Delta(\partial \Delta_n)=\Psi(\partial K)=\mathbb{S}^{n-1}.
% \]
% \end{proof}

% \begin{remark}
% We collect several definitions related to convex sets in $\mathbb{R}^n$ below.
% \end{remark}

% \begin{definition}[Relative interior]
% For $K\subset\R^n$, with affine hull $\aff(K)$, the \emph{relative interior} is
% \[
% \ri(K):=\Int_{\aff(K)}(K).
% \]
% If $\aff(K)=\mathbb{R}^n$ then $\ri(K)=\Int(K)$; otherwise $\ri(K)$ is the interior of $K$ taken in its own affine hull $\aff(K)$.
% \end{definition}

% \begin{proposition}[Convex characterization ]
% If $K\subset\R^n$ is convex and $x\in K$, the following are equivalent:
% \begin{enumerate}
% \item $x\in\ri(K)$.
% \item There exists $\varepsilon>0$ with $x+\varepsilon(B\cap V)\subset K$, where
% $B=\{y:\|y\|\le 1\}$ and $V=\Span(K-K)$.
% \item No supporting hyperplane of $K$ is active at $x$ (i.e.\ $x$ lies on no proper face).
% \end{enumerate}
% \end{proposition}

% \begin{lemma}
% If $K=\conv\{v_0,\dots,v_m\}$ and $x=\displaystyle \sum_{i=0}^m \lambda_i v_i$ with all
% $\lambda_i>0$, $\displaystyle\sum_i\lambda_i=1$, then $x\in\ri(K)$ (hence not on any proper face).
% \end{lemma}
% \begin{proof}
% Suppose, for a contradiction, that $x\notin\ri(K)$. Then there exists a supporting
% hyperplane $H=\{y\in\aff(K): a\cdot y=b\}$ of $K$ such that $K\subset\{a\cdot y\le b\}$
% and $x\in H$ (i.e.\ the supporting inequality is \emph{active} at $x$).

% The corresponding face is $F:=K\cap H$, which is a proper face of $K$.

% Because $K=\conv\{v_0,\dots,v_m\}$, the face $F$ is the convex hull of those
% vertices lying on $H$:
% \[
% I:=\{\,i:\ a\cdot v_i=b\,\},\qquad F=\conv\{v_i:\ i\in I\}.
% \]
% Since $F$ is proper, $I\subsetneq\{0,\dots,m\}$,
% so there exists $j$ with $a\cdot v_j<b$.

% On the other hand
% \[
% a\cdot x=\sum_{i=0}^m \lambda_i\, a\cdot v_i
% =\sum_{i\in I}\lambda_i\, b \;+\; \sum_{i\notin I}\lambda_i\, a\cdot v_i
% < \Bigl(\sum_{i=0}^m \lambda_i\Bigr)b \;=\; b,
% \]
% since at least one index with $a\cdot v_i<b$ has $\lambda_i>0$.
% This contradicts $x\in H$ (which requires $a\cdot x=b$).

% Hence no supporting hyperplane of $K$ is active at $x$, i.e.\ $x$ lies on no
% proper face. By the standard characterization of relative interior in convex
% sets, this is equivalent to $x\in\ri(K)$.
% \end{proof}

% \begin{definition}[Convex hull]
% For $S\subset\R^n$,
% \[
% \conv(S)=\bigcap\{\,C\supset S:\ C\text{ convex}\,\}
% =\left\{\sum_{i=1}^m \lambda_i x_i:\ x_i\in S,\ \lambda_i\ge 0,\ \sum_i\lambda_i=1\right\}.
% \]
% \end{definition}

% \begin{definition}[Finitely generated cone and basic types]
% Given $v_1,\dots,v_m\in\R^n$,
% \[
% \cone(v_1,\dots,v_m):=\Bigl\{\sum_{i=1}^m \alpha_i v_i:\ \alpha_i\ge 0\Bigr\}.
% \]
% For a cone $C$, set $\lin(C):=\{x:\ x,-x\in C\}$ and $\Span(C):=\Span\{C\}$.
% Its dimension is $\dim C:=\dim(\Span(C))$. The cone is \emph{pointed} if
% $\lin(C)=\{0\}$, and \emph{simplicial} if its generators spanning $\Span(C)$
% are linearly independent (then $\dim C$ equals their number).
% \end{definition}

% \begin{definition}[Dimension of a finitely generated cone]
% If $C=\cone(v_1,\dots,v_m)$ and $V=[v_1\ \cdots\ v_m]$, then
% \[
% \dim C:=\dim\Span\{v_1,\dots,v_m\}=\rank(V).
% \]
% \end{definition}

\section{Standard Maps between Cells and Spheres}\label{sec:cells-spheres}

\begin{definition}\label{def:std-ball-sphere-cell}
\begin{enumerate}
    \item The standard $n$-ball
    \[\B^n = \{x \in \R^n: \norm{x} \leq 1\}.\]
    
    \item The standard $n$-sphere
    \[\mathbb{S}^n= \partial(\B^{n+1}) = \{x \in \R^{n+1}: \norm{x} =1\}.\]

    \item The standard $n$-cell is a topological space that is homeomorphic to the open unit ball 
    \[\Int(\B^n) = \{x \in \R^n: \norm{x} < 1\}.\]
\end{enumerate}
\end{definition}

\begin{proposition}[Stereographic projection]\label{prop:stereo}
    For $n \geq 0$, denote 
    \[N=(0,\ldots,0,1),~S=(0,\ldots,0,-1)\in \mathbb{S}^n\]
    called the north pole and the south pole of $\mathbb{S}^n$, respectively. Then $\mathbb{S}^n\setminus\{S\}$ and $\mathbb{S}^n\setminus\{N\}$ are homeomorphic to $\R^n$.
\end{proposition}

\begin{proof}
    Consider the continuous maps
    \[
        f:\mathbb{S}^n\setminus\{N\}\to \R^n,~
        x=(a,t) \in \mathbb{S}^n\setminus\{N\} \cap \R^n \times \R\mapsto \dfrac{1}{1-t}a,
    \]
    \[
        g:\R^n\to \mathbb{S}^n\setminus\{N\}, \quad 
        y \mapsto (\phi(y)\cdot y,1-\phi(y)),
    \]
    where $\phi(y)=\dfrac{2}{1+\|y\|^2}$.

    For every $x=(a,t)\in \mathbb{S}^n$, we compute
    \[
        \phi(f(x))=\phi\left(\dfrac{1}{1-t}a\right)=\dfrac{2}{1+\dfrac{1}{(1-t)^2}\norm{a}^2} = \dfrac{(1-t)^2}{1-2t+t^2+\norm{a}^2}=\dfrac{(1-t)^2}{2-2t}
        =1-t,
    \]
    hence
        \[(g\circ f)(x)
        =g\left(\dfrac{1}{1-t}a\right)
        =\left(\phi(f(x))\cdot f(x),1-\phi(f(x))\right)
        =(a,t)=x.\]

    For every $y\in \R^n$, we have
        \[(f\circ g)(y)
        =f(\phi(y)\cdot y,1-\phi(y))
        = \dfrac{1}{1-(1-\phi(y))}\left(\phi(y)\cdot y\right)
        =y.\]

    Therefore $f$ is a continuous bijection with continuous inverse $f^{-1}=g$, i.e.\ $f$ is a homeomorphism between $\mathbb{S}^n\setminus\{N\}$ and $\R^n$.

    On the other hand, there is also a homeomorphism.
    \[
        \mathbb{S}^n\setminus\{N\}\to \mathbb{S}^n\setminus\{S\}, \quad 
        (a,t)\mapsto (a,-t).
    \]
\end{proof}

\begin{proposition}\label{prop:open-ball-Rn}
    The open unit ball $\Int(\B^n)$ 
    is homeomorphic to $\R^n$.
\end{proposition}

\begin{proof}
    We have constructed explicit continuous maps
    \[
        \varphi:\Int(\B^n)\to \R^n, \quad x \mapsto \dfrac{x}{1-\norm{x}}, 
    \]
    and
    \[
        \psi:\R^n\to \Int(\B^n), \quad y \mapsto \dfrac{y}{1+\|y\|}.
    \]
    Both $\varphi$ and $\psi$ are continuous and
    $\psi\circ \varphi =\Id_{\Int(\B^n)},~ \varphi\circ \psi=\Id_{\R^n}$. 
    
    Hence $\varphi$ is a homeomorphism
    \[
        \Int(\B^n) \approx \R^n
    \]
\end{proof}

\begin{remark}\label{rmk:pi-assembly}
    \begin{enumerate}
        \item From \Cref{prop:open-ball-Rn,prop:stereo}, we obtain the homeomorphisms
        \[\Int(\B^n) \xlongrightarrow[\varphi]{\approx}\R^n \xlongrightarrow[g]{\approx} \mathbb{S}^n\setminus\{N\}.\]
    
        \item Since 
        \[\B^n = \Int(\B^n) \sqcup \partial(\B^n)=\Int(\B^n) \sqcup \mathbb{S}^{n-1}\] and 
        \[\mathbb{S}^n = \mathbb{S}^n\setminus\{N\} \sqcup \{N\}.\]
        Naturally, we want to define a map
        \[\pi: \B^n \to \mathbb{S}^n, \]
        such that 
        \[\pi(\Int(\B^n)) = \mathbb{S}^{n}\setminus \{N\}= (g \circ \varphi)(\Int(\B^n)),\quad  \pi(\mathbb{S}^{n-1}) = \{N\}.\]
        With $x \in \Int(\B^n),~y = \varphi(x) = \dfrac{x}{1-\norm{x}}$, so 
            \[\phi(y)=\dfrac{2}{1+\norm{y}^2}
            =\dfrac{2}{1+\dfrac{\norm{x}^2}{(1-\norm{x})^2}} 
            = \dfrac{2(1-\norm{x})^2}{\norm{x}^2 + (1-\norm{x})^2}\]
        
       \[(g\circ \varphi)(x)
            = \left(\phi(y)\cdot y, 1-\phi(y)\right)
            =\dfrac{1}{\norm{x}^2+(1-\norm{x})^2}\left(2(1-\norm{x})x, 2\norm{x}-1\right) \in \mathbb{S}^n.\]
        \end{enumerate}
    Moreover, with $z = (a,t) \in \mathbb{S}^n\setminus\{N\} \cap \R^n \times \R$
        \[(g\circ \varphi)^{-1}(z)
        =(\psi\circ f)(z)
        =\psi\left(\dfrac{1}{1-t}a\right)
        = \dfrac{a}{(1-t)+\norm{a}} \in \Int(\B^n).\]
\end{remark}

Hence, we have a definition, also a proposition following:
\begin{proposition}\label{prop:quotient-pi}
    The map 
    \[\pi:(\B^n, \mathbb{S}^{n-1})\to (\mathbb{S}^n,N),~x \mapsto \dfrac{1}{\norm{x}^2+(1-\norm{x})^2}\left(2(1-\norm{x})x, 2\norm{x}-1\right)\]
    induces a homeomorphism \[\overline{\pi}:\B^n/\mathbb{S}^{n-1} \approx \mathbb{S}^n.\]
    
    Moreover, $\pi \mid_{\Int(\B^n)}: \Int(\B^n) \approx \mathbb{S}^n\setminus\{N\}$ is a homeomorphism and the inverse of $\pi \mid_{\Int(\B^n)}$ is
    \[\rho = \psi \circ f: \mathbb{S}^n\setminus\{N\} \to \Int(\B^n),~\]
\end{proposition}

\begin{remark}\label{rmk:quotient-remark}
    $\B^n/\mathbb{S}^{n-1} = \B^n/\sim$, where $x\sim y$ if $x = y$ or $x,y\in \mathbb{S}^{n-1}$. 
\end{remark}

Now, we will show that simplices, cubes and their products are homeomorphic with balls.

\begin{lemma}[Normalization]\label{lem:normalization}
    Let $B\subset \R^n$ be a closed $n$--ball of radius $r>0$ centered at $c\in\R^n$. 
    Then there exists an affine homeomorphism 
    \[
        T:\R^n \to \R^n
    \]
    such that 
    \[
        T(B) = \B^n \quad \text{and} \quad T(\partial B) = \partial \B^n.
    \]
    Consequently, 
    \[
        (B,\partial B)\approx (\B^n,\partial \B^n).
    \]
\end{lemma}

\begin{proof}
    Define 
    \[
        T:\R^n \to \R^n, \qquad T(x) = \dfrac{1}{r}(x-c).
    \]
    This is the composition of a translation by $-c$ and a dilation by factor $1/r$, 
    hence a homeomorphism of $\R^n$ onto itself. 

    For $x\in B$ we have $\|x-c\|\le r$, so 
    \[
        \|T(x)\| = \left\|\dfrac{x-c}{r}\right\|\le 1,
    \]
    i.e.\ $T(x)\in \B^n$. Conversely, if $y\in \B^n$, then $\|y\|\le 1$ and 
    $x = r y + c\in B$, so $T(B)=\B^n$. The same argument shows 
    $T(\partial B)=\partial \B^n$. 

    Therefore $T$ restricts to a homeomorphism of pairs 
    $(B,\partial B)\approx (\B^n,\partial \B^n)$.
\end{proof}

\begin{proposition}\label{prop:convex-standard}
    If $K$ is a compact convex subset of $\R^n$ which contains an $n$-ball $B$ then there is a standard homeomorphism $(B,\partial(B)) \approx (K,\partial(K))$. 
\end{proposition}
\begin{proof}
    Using the \Cref{lem:normalization}, we only need prove $(\B^n,\partial(\B^n)) \approx (K,\partial(K))$.

    For each $u\in\mathbb{S}^{n-1}$, put
    \[
        \lambda(u) :=\sup\{ t\ge 0 : tu\in K\},
    \]
    then 
    \[1 \leq \lambda(u) < \infty.\] 
    Indeed, the compactness of $K$ implies there exists  $R>0$ such that $\norm{x}\leq R$ for all $x\in K$, hence $\lambda(u) <\infty$. Moreover, $\lambda(u)\ge 1$, since
    \[[0,1]\cdot u = \{tu: t\in [0,1]\}\subset \B^n \subset K.\]

    By the definition, $\lambda(u)u$ is the \textit{furthest point of} $K$ along the ray $\{tu:t\geq 0\}$ and cannot be outside $K$, because all points $tu$ with $t < \lambda(u)$ are in $K$ and $K$ is closed (compact). 
    
    If $\lambda(u)u \in \Int(K)$, there exists $\varepsilon > 0$ such that $(\lambda(u) + \varepsilon)u \in K$, a contradiction of the definition of $\lambda(u)$. So $\lambda(u)u \in \partial K$ and is the unique intersection of the ray $\{t u:t\ge 0\}$ with $\partial K$.

    Now, we prove that $\lambda$ is continuous. 
    
    Let $u_k\to u$ in $\mathbb{S}^{n-1}$ and put $t_k=\lambda(u_k)$.
    Since $K$ is compact, $(t_k)$ is bounded; pass to a convergent subsequence
    $t_k\to t$. Because $K$ is closed, $t u=\lim t_k u_k\in K$, hence $t\le \lambda(u)$.
    For the reverse inequality, fix $\varepsilon>0$. Then
    $(\lambda(u)-\varepsilon)u \in \Int(K)$; by openness there is a neighborhood
    $U$ of $u$ with $(\lambda(u)-\varepsilon) v\in K$ for all $v\in U$, hence
    $\lambda(v)\ge \lambda(u)-\varepsilon$ for $v\in U$. Taking $v=u_k$ close to $u$ gives
    $\liminf t_k\ge \lambda(u)-\varepsilon$. Since $\varepsilon>0$ is arbitrary,
    $\liminf t_k\ge \lambda(u)$. Thus $t=\lambda(u)$, so $\lambda$ is continuous.
    
    Finally, we contruct a homeomorphim between $\B^n$ and $K$. 
    
    For all $x\in\B^n$, can be written uniquely as $x=ru$ with $u\in\mathbb{S}^{n-1}$ and
    $r\in[0,1]$. Define
    \[
        F:\B^n\longrightarrow K,\qquad x= r u \mapsto r\lambda(u)u.
    \]
    This is well-defined and continuous (product of continuous functions on the
    polar coordinates). It is bijective with inverse
    \[
        F^{-1}:K\longrightarrow \B^n,\qquad 
        F^{-1}(s v)= \dfrac{s}{\lambda(v)}\,v\quad\bigl(v\in\mathbb{S}^{n-1},\ s\in[0,\lambda(v)]\bigr),
    \]
    which is also continuous by continuity of $\lambda$. Hence $F$ is a homeomorphism.
    
    In particular, if $r=1$ then $F(u)=\lambda(u)u\in\partial K$; if $r<1$ then $F(r u)$ lies on the
    segment between $0$ and $\lambda(u)u$ and hence in $\Int(K)$.
    Therefore $F(\partial\B^n)=\partial K$ and $F(\Int\,\B^n)=\Int(K)$. Thus $F:(\B^n,\partial\B^n)\approx (K,\partial K)$.
\end{proof}

\begin{corollary}\label{cor:cube-ball}
    The cube \(I^n=[-1,1]^n\subset\R^n\) is homeomorphic to the (closed) unit ball
\(\B^n=\{x\in\R^n:\|x\|\le 1\}\). Moreover, this homeomorphism sends
\(\partial I^n\) onto \(\mathbb{S}^{n-1}=\partial(\B^n)\)
\end{corollary}
\begin{proof}
    Since $\norm{x} \leq 1$ implies $|x_j|\le 1$ for each coordinate $x_j$. Thus $I^n$ is a compact convex set of $\R^n$ containing the unit $n$–ball. 
    
    \Cref{prop:convex-standard} applies with $K=I^n$ and $B=\B^n$, yielding a standard homeomorphism
    \[
        F:(\B^n,\mathbb{S}^{n-1})\xrightarrow{\ \approx\ }(I^n,\partial I^n),~
        ru \mapsto r\,\lambda(u)\,u 
        \qquad (u \in \mathbb{S}^{n-1}, r \in [0,1]).
    \]
Fix $u\in\mathbb{S}^{n-1}$, for $t\ge 0$,
\[
tu\in I^n 
\iff t u_i \in I,~\forall i
\iff |t u_i|\le 1,~\forall i
\iff t\le \dfrac{1}{|u_i|}\ \text{whenever }u_i\neq 0.
\]
Hence the set 
\[\{t\ge 0: tu\in I^n\}=
\Bigl[\,0,\ \min_{1\le i\le n,\,u_i\neq 0}\dfrac{1}{|u_i|}\,\Bigr],
\]
so its supremum is
\[
\lambda(u)=\min_{1\le i\le n,\,u_i\neq 0}\dfrac{1}{|u_i|}
=\dfrac{1}{\max_{1\le i\le n}|u_i|}
=\dfrac{1}{\|u\|_\infty}.
\]
Taking $t=\dfrac{1}{\|u\|_\infty}$ gives 
$|t u_i|=\dfrac{|u_i|}{\|u\|_\infty}\le 1$ for all $i$, with equality for at least one index attaining the maximum of $|u_i|$. Thus $tu\in\partial I^n$, so the supremum is attained and the intersection with $\partial I^n$ is exactly at $u/\|u\|_\infty$.

Consequently, its inverse \(F^{-1}:I^n\to\B^n\) is
\[
F^{-1}(x)=
\begin{cases}
0, & x=0,\\[4pt]
\dfrac{\norm{x}_\infty}{\norm{x}}\,x, & x\neq 0,
\end{cases}
\]
which is continuous on \(I^n\) and bijective with continuous inverse.

Finally, if \(x\in\partial I^n\) then \(\norm{x}_\infty=1\), so
\(\|F^{-1}(x)\|=\norm{x}_\infty=1\) and hence \(F^{-1}(x)\in\mathbb{S}^{n-1}\).  
If \(\norm{x}_\infty<1\) (i.e. \(x\in \Int I^n\)), then
\(\|F^{-1}\|=\norm{x}_\infty<1\), so \(F^{-1}(x)\in \Int\B^n\).
Therefore \(F^{-1}(\partial I^n)=\mathbb{S}^{n-1}\) and
\(F^{-1}(\Int I^n)=\Int\B^n\), completing the proof.
\end{proof}

\begin{corollary}\label{cor:simplex-ball}
    Let the standard $n$–simplex be
    \[
        \Delta_n=\Bigl\{(t_0,\dots,t_n)\in\mathbb{R}^{n+1}_{\ge 0}\ :\ \sum_{i=0}^n t_i=1\Bigr\},
    \]
    and let $\partial \Delta_n=\{t\in\Delta_n:\ t_i=0\ \text{for some }i\}$ be its boundary.
    
    There exists a linear embedding $i:\Delta_n\hookrightarrow\mathbb{R}^n$ such that $K:=\imath (\Delta_n)$ is a compact convex set containing a Euclidean $n$–ball. Consequently, by radial projection there is a \emph{standard} homeomorphism
    \[
        \Phi_\Delta:(\Delta_n,\partial \Delta_n)\;\xrightarrow{\ \cong\ }\;(\B^{\,n},\mathbb{S}^{n-1}).
    \]
\end{corollary}

\begin{proof}
    Define a linear embedding $\imath:\Delta_n\to\mathbb{R}^n$ by sending the vertices as follows:
for $0\le j\le n-1$ let $v_j=e_j$ (the $j$th standard basis vector of $\mathbb{R}^n$) and let
\[
   v_n=-\sum_{j=0}^{n-1} e_j.
\]
On barycentric coordinates $t=(t_0,\dots,t_n)\in\Delta_n$ we put
\[
   \imath(t)=\sum_{j=0}^n t_j v_j\in\mathbb{R}^n.
\]
Then $K:=\imath(\Delta_n)=\operatorname{conv}\{v_0,\dots,v_n\}$ is a compact convex polytope in $\mathbb{R}^n$. We have two claims following
\begin{itemize}
    \item $0\in\Int(K)$.
    
Indeed,
\[
   \frac{1}{n+1}\sum_{j=0}^n v_j
   =\frac{1}{n+1}\Bigl(\sum_{j=0}^{n-1} e_j - \sum_{j=0}^{n-1} e_j\Bigr)=0,
\]
so $0$ is a convex combination of the vertices with all coefficients equal to $1/(n+1)>0$; hence $0$ lies in the (relative) interior of $K$.

\item \textit{$K$ contains a Euclidean $n$–ball around $0$.}

Since $K$ is compact and $0\notin\partial K$, the continuous function $x\mapsto\|x\|$ attains a positive minimum on the compact set $\partial K$:
\[
   r:=\min\{\|x\|:\ x\in\partial K\}>0.
\]
Therefore the closed ball $\overline B(0,r)\subset K$.

By \Cref{prop:convex-standard} applied to $K$ (with the ball $\overline B(0,r)$ found above), there is a standard homeomorphism
\[
   \Psi:(K,\partial K)\xrightarrow{\approx}( \B^{\,n},\mathbb{S}^{n-1})
\]
that sends each $y\in\partial K$ to $y/\|y\|$ and extends linearly on the radial segments from $0$.
\end{itemize}
Finally, set $\Phi_\Delta:=\Psi\circ \imath:\Delta_n \to \B^{n}$. Since $i$ is a homeomorphism from $\Delta_n$ onto $K$ and $\Psi$ is a homeomorphism from $K$ onto $\B^{\,n}$, $\Phi_\Delta$ is a homeomorphism.

It remains to check the boundary correspondence. 

A point $t\in\Delta_n$ lies in $\partial \Delta_n$ iff at least one barycentric coordinate is $0$, say $t_j=0$; equivalently, $i(t)$ lies in the facet $\operatorname{conv}\{v_0,\dots,\widehat{v_j},\dots,v_n\}\subset\partial K$. Thus $i(\partial \Delta_n)=\partial K$, and hence
\[
   \Phi_\Delta(\partial \Delta_n)=\Psi(\partial K)=\mathbb{S}^{n-1}.
\]
\end{proof}

\begin{remark}\label{rmk:convex-collect}
We collect several definitions related to convex sets in $\mathbb{R}^n$ below.
\end{remark}

\begin{definition}[Relative interior]\label{def:ri}
For $K\subset\R^n$, with affine hull $\aff(K)$, the \emph{relative interior} is
\[
\ri(K):=\Int_{\aff(K)}(K).
\]
If $\aff(K)=\mathbb{R}^n$ then $\ri(K)=\Int(K)$; otherwise $\ri(K)$ is the interior of $K$ taken in its own affine hull $\aff(K)$.
\end{definition}

\begin{proposition}[Convex characterization]\label{prop:ri-equiv}
If $K\subset\R^n$ is convex and $x\in K$, the following are equivalent:
\begin{enumerate}
\item $x\in\ri(K)$.
\item There exists $\varepsilon>0$ with $x+\varepsilon(B\cap V)\subset K$, where
$B=\{y:\|y\|\le 1\}$ and $V=\Span(K-K)$.
\item No supporting hyperplane of $K$ is active at $x$ (i.e.\ $x$ lies on no proper face).
\end{enumerate}
\end{proposition}

\begin{lemma}\label{lem:ri-barycentric}
If $K=\conv\{v_0,\dots,v_m\}$ and $x=\displaystyle \sum_{i=0}^m \lambda_i v_i$ with all
$\lambda_i>0$, $\displaystyle\sum_i\lambda_i=1$, then $x\in\ri(K)$ (hence not on any proper face).
\end{lemma}
\begin{proof}
Suppose, for a contradiction, that $x\notin\ri(K)$. Then there exists a supporting
hyperplane $H=\{y\in\aff(K): a\cdot y=b\}$ of $K$ such that $K\subset\{a\cdot y\le b\}$
and $x\in H$ (i.e.\ the supporting inequality is \emph{active} at $x$).

The corresponding face is $F:=K\cap H$, which is a proper face of $K$.

Because $K=\conv\{v_0,\dots,v_m\}$, the face $F$ is the convex hull of those
vertices lying on $H$:
\[
I:=\{\,i:\ a\cdot v_i=b\,\},\qquad F=\conv\{v_i:\ i\in I\}.
\]
Since $F$ is proper, $I\subsetneq\{0,\dots,m\}$,
so there exists $j$ with $a\cdot v_j<b$.

On the other hand
\[
a\cdot x=\sum_{i=0}^m \lambda_i\, a\cdot v_i
=\sum_{i\in I}\lambda_i\, b \;+\; \sum_{i\notin I}\lambda_i\, a\cdot v_i
< \Bigl(\sum_{i=0}^m \lambda_i\Bigr)b \;=\; b,
\]
since at least one index with $a\cdot v_i<b$ has $\lambda_i>0$.
This contradicts $x\in H$ (which requires $a\cdot x=b$).

Hence no supporting hyperplane of $K$ is active at $x$, i.e.\ $x$ lies on no
proper face. By the standard characterization of relative interior in convex
sets, this is equivalent to $x\in\ri(K)$.
\end{proof}

\begin{definition}[Convex hull]\label{def:conv-hull}
For $S\subset\R^n$,
\[
\conv(S)=\bigcap_{C \supset S}\{C:\ C\text{ convex}\,\}
=\left\{\sum_{i=1}^m \lambda_i x_i:\ x_i\in S,\ \lambda_i\ge 0,\ \sum_i\lambda_i=1\right\}.
\]
\end{definition}

\begin{definition}[Finitely generated cone and basic types]\label{def:cone}
Given $v_1,\dots,v_m\in\R^n$,
\[
\cone(v_1,\dots,v_m):=\Bigl\{\sum_{i=1}^m \alpha_i v_i:\ \alpha_i\ge 0\Bigr\}.
\]
For a cone $C$, set $\lin(C):=\{x:\ x,-x\in C\}$ and $\Span(C):=\Span\{C\}$.
Its dimension is $\dim C:=\dim(\Span(C))$. The cone is \emph{pointed} if
$\lin(C)=\{0\}$, and \emph{simplicial} if its generators spanning $\Span(C)$
are linearly independent (then $\dim C$ equals their number).
\end{definition}

\begin{definition}[Dimension of a finitely generated cone]\label{def:cone-dim}
If $C=\cone(v_1,\dots,v_m)$ and $V=[v_1\ \cdots\ v_m]$, then
\[
\dim C:=\dim\Span\{v_1,\dots,v_m\}=\rank(V).
\]
\end{definition}
